\chapter*{ABSTRACT}
\addcontentsline{toc}{chapter}{ABSTRACT}

\begin{center}

    \textbf{Prediction-Conditioned RAG: Integrating Glucose Level Prediction and Document-Based Clinical Recommendations for Diabetes Management} \\
    \vspace{0.4cm}
    by \\
    \vspace{0.4cm}
    Daffari Adiyatma\\
    18222003\\
\end{center}


\begin{spacing}{1.0}
Clinical decision support systems for diabetes that employ glucose prediction or
\emph{Retrieval-Augmented Generation} (RAG) generally retrieve knowledge based on the patient's
current condition, so their recommendations are reactive. No prior study has been found that
conditions the retrieval process on the predicted glucose level. This final project proposes
prediction-conditioned RAG, a mechanism that composes the retrieval query from the predicted
glucose level so that the retrieved documents and the resulting recommendations are anticipatory,
within a workflow whose output is always verified by a physician. Glucose prediction is performed
by a \emph{Gradient Boosting} model with physiologically derived features, while clinical knowledge
is retrieved from 2,233 guideline chunks through lexical retrieval. Evaluation on the OhioT1DM
dataset uses six-fold cross-patient cross-validation. At the $+30$-minute horizon, the model
achieves an RMSE of 20.72~mg/dL with 94.69\% of points in zones A+B of the \emph{Clarke Error
Grid}, comparable to \emph{Random Forest} and LSTM; in the self-monitoring scenario without CGM
(horizon of about 240 minutes) the error rises to 71.60~mg/dL but remains below the persistence
\emph{baseline} (95.31~mg/dL). In divergent cases, where the future condition differs from the
current one, prediction-conditioned RAG consistently outperforms standard RAG across all six
\emph{folds} (\emph{Mean Reciprocal Rank} 0.190 to 0.284; $p=0.031$) and on 120 real cases
(0.186 to 0.321; $p<0.001$). The \emph{oracle} arm reaches an MRR of 0.778, so most of the
remaining gap stems from prediction accuracy rather than from the retrieval mechanism. The same
limitation appears in hypoglycemia sensitivity, namely 17.3\% with an upper threshold on the
regression output and 67.2\% with a direct condition classifier. All local computation runs
without a GPU (median 0.339~s, memory 342~MB), and every reference can be traced to the printed
page of its guideline. A preliminary evaluation by twelve healthcare professionals rated the
recommendations as consistent with the guidelines under physician verification, with a
\emph{System Usability Scale} score of 55.21 indicating that the system's usability still needs
improvement.

\textbf{Keywords:} prediction-conditioned RAG, glucose level prediction, \emph{Gradient Boosting}, clinical decision support system, diabetes.
\end{spacing}
